\section{Christ the Judge: The Father Has Given All Judgment to the Son}

On every Sunday and solemnity of the Christian year, the faithful rise to profess the Creed and look beyond history toward a single return: ``from thence he shall come to judge the living and the dead''---\textit{inde venturus est judicare vivos et mortuos}---or, in the Nicene-Constantinopolitan confession, ``and he shall come again with glory to judge the living and the dead, whose kingdom shall have no end.''\endnote{Apostles' Creed, Article VII; Nicene-Constantinopolitan Creed.} At the Fourth Lateran Council in 1215, the Church gathered this faith into one solemn sentence: the incarnate Son of God, who suffered on the Cross, descended in soul, rose in flesh, and ascended in both, is ``to come at the end of the world, to judge the living and the dead, and to render to every one according to his works, both to the reprobate and to the elect''---\textit{venturus in fine saeculi, judicaturus vivos et mortuos, et redditurus singulis secundum opera sua, tam reprobis quam electis}.\endnote{Fourth Lateran Council (1215), Constitution 1, \textit{Firmiter credimus} (DS 801).} Judgment is neither an impersonal law of nature nor a mechanical tally kept by fate; it is a personal encounter with Jesus Christ.

The \textit{Roman Catechism}, opening its exposition of the seventh article of the Apostles' Creed, sets Christ's judiciary office beside His Passion and His heavenly intercession:

\begin{quote}
Jesus Christ is invested with three eminent offices and functions, those of Redeemer, Patron, and Judge. But as, in the preceding Articles, we have shown that the human race was redeemed by his passion and death; and as, by his ascension into heaven, it is manifest that he has undertaken the perpetual advocacy and patronage of our cause; it next follows, that, in this Article, we set forth his character as judge.\endnote{\textit{Catechism of the Council of Trent}, Part I, Creed Article VII, trans.\ J.\ Donovan (Baltimore: Lucas Brothers, 1829), p.\ 61.}
\end{quote}

To sever the Judge from the Redeemer and the Advocate is to misread the Gospel. The One before whom every conscience must stand is the very Son who took flesh from the Virgin Mary, poured out His blood upon Calvary, and stands before the Father as our Patron.

Sacred Scripture announces this commission directly from the lips of Christ. In the fifth chapter of St.\ John's Gospel, our Lord binds His eternal Sonship to His assumed humanity:

\begin{quote}
For neither does the Father judge any man: but hath given all judgment to the Son. That all men may honour the Son, as they honour the Father. He who honoureth not the Son honoureth not the Father who hath sent him. \dots\ For as the Father hath life in himself, so he hath given to the Son also to have life in himself: And he hath given him power to do judgment, because he is the Son of man. Wonder not at this: for the hour cometh wherein all that are in the graves shall hear the voice of the Son of God. And they that have done good things shall come forth unto the resurrection of life: but they that have done evil, unto the resurrection of judgment. (John 5:22--23, 26--29)\endnote{Scripture quotations throughout follow the Challoner Douay-Rheims Bible, corroborated against the Clementine Vulgate; Psalm and book numbering follows the Vulgate.}
\end{quote}

The apostolic preaching repeats this charge as a pillar of the Gospel. In the house of Cornelius, St.\ Peter testifies ``that it is he who was appointed by God to be judge of the living and of the dead'' (Acts 10:42). At the Areopagus, St.\ Paul proclaims that God ``hath appointed a day wherein he will judge the world in equity, by the man whom he hath appointed: giving faith to all, by raising him up from the dead'' (Acts 17:31). To the Corinthians, the Apostle names the tribunal itself: ``For we must all be manifested before the judgment seat of Christ, that every one may receive the proper things of the body, according as he hath done, whether it be good or evil'' (2 Cor 5:10).

In the Old Testament, judgment is attributed to the Lord God Himself---the Ancient of days whose throne is flames of fire (Dan 7:9--10). St.\ Augustine demonstrates in \textit{The City of God} that whenever the prophets announce God's coming to judge, ``we must gather that Christ is meant; for though the Father will judge, He will judge by the coming of the Son.''\endnote{St.\ Augustine, \textit{The City of God} XX.30, trans.\ Marcus Dods (Edinburgh: T.\ \& T.\ Clark, 1871).} St.\ Cyril of Jerusalem guards the same Trinitarian unity: the Father commits judgment to the Son ``not divesting Himself of His power, but judging through the Son; the Son therefore judgeth by the will of the Father; for the wills of the Father and of the Son are not different, but one and the same.''\endnote{St.\ Cyril of Jerusalem, \textit{Catechetical Lectures} XV.25, in \textit{Nicene and Post-Nicene Fathers}, Series II, vol.\ 7.}

St.\ Thomas Aquinas unfolds Christ's judiciary power systematically in Question 59 of the \textit{Tertia Pars}.\endnote{St.\ Thomas Aquinas, \textit{Summa Theologiae} III, q.\ 59, aa.\ 1--6, trans.\ Fathers of the English Dominican Province.} Although judiciary authority belongs to the entire Blessed Trinity in its undivided essence, it is specially appropriated to the Son. Three things are required for judgment: power to coerce subjects, upright zeal from love of justice, and wisdom as the rule of judgment itself. Because the Son is ``Wisdom begotten, and Truth proceeding from the Father, and His perfect Image,'' whatever the Father judges, He judges through His Word. In Daniel's vision (Dan 7:9--14), the Ancient of days signifies that judicial authority originates in the Father, while the Son of man receives power, honor, and an everlasting kingdom from Him.\endnote{Aquinas, \textit{ST} III, q.\ 59, a.\ 1, corpus and ad 1--2.}

Yet John 5:27 adds ``because he is the Son of man.'' Why does universal judgment belong to Christ also in His human nature? Heeding St.\ John Chrysostom's caution that this authority does not spring merely from common human nature---else every man would be universal judge---St.\ Thomas shows that it belongs to Christ as man by the \textit{grace of Headship}, because even in His humanity He is Head of the Church and of all mankind, under whose feet God has placed all things.\endnote{Aquinas, \textit{ST} III, q.\ 59, a.\ 2, corpus.}

From this Headship, St.\ Thomas draws three fitting reasons. First, kinship: God judges men through the Man Christ ``that His judgment may be sweeter to men,'' echoing Hebrews 4:15--16 on our compassionate High Priest tempted in all things without sin. Second, the bodily resurrection: as St.\ Augustine observes, God raises dead souls through the Son of God and raises dead bodies through the Son of man. Third, visibility to all who are judged: because the uncreated divine essence cannot be beheld without supreme joy---as the posthumous \textit{Supplement} proves from Aquinas's \textit{Commentary on the Sentences}, ``it is impossible for the Godhead to be seen without joy''\endnote{\textit{Supplementum Tertiae Partis} (compiled from Aquinas's \textit{Scriptum super Sententiis}, Book IV), q.\ 90, a.\ 3, corpus, trans.\ Fathers of the English Dominican Province (1922).}---St.\ Augustine and St.\ Thomas conclude that ``the form of a servant should be shown in the judgment to both good and wicked, while the form of God should be kept for the good alone.''\endnote{Aquinas, \textit{ST} III, q.\ 59, a.\ 2, corpus (citing St.\ Augustine, \textit{Sermon} 127).} Thus, as the \textit{Roman Catechism} notes, men ``may see their judge with their eyes, and hear him with their ears, and thus learn their final doom, through the medium of the senses.''\endnote{\textit{Catechism of the Council of Trent}, Part I, Creed Article VII, trans.\ Donovan, p.\ 63.}

Christ also \textit{merited} this judiciary dignity by His obedience unto death. While it is already His by hypostatic union, Headship, and fullness of grace, ``in accordance with the Divine justice, He should be judge who fought for God's justice, and conquered, and was unjustly condemned.''\endnote{Aquinas, \textit{ST} III, q.\ 59, a.\ 3, corpus.} The \textit{Roman Catechism} puts it starkly: ``Most just is it also, that he who was most iniquitously condemned by the judgment of men, should himself be, afterwards seen by all men sitting in judgment on all.''\endnote{\textit{Catechism of the Council of Trent}, Part I, Creed Article VII, trans.\ Donovan, p.\ 63.} St.\ Augustine compresses the reversal into a single aphorism: ``The Judge shall sit, who stood before a judge; He shall condemn the truly wicked, who Himself was falsely reputed wicked.''\endnote{Quoted in Aquinas, \textit{ST} III, q.\ 59, a.\ 3, sed contra (citing St.\ Augustine, \textit{Sermon} 127).} And the \textit{Supplement} adds that Christ is Lord of the dead and the living not only by Creation as God, ``but also by reason of the Redemption, which pertains to Him in respect of His human nature'' (Rom 14:9), so that He who opened the kingdom by His blood fittingly presides over admission to it.\endnote{\textit{Supplementum}, q.\ 90, a.\ 1, corpus.}

Nothing in creation falls outside this jurisdiction. Though our Lord declined to settle a temporal inheritance dispute on earth (Luke 12:13--14), St.\ Thomas demonstrates that all human affairs lie under Christ's tribunal because every human act is ordered to eternal salvation; even the secrets of the heart are searched by Christ's human soul through the light overflowing from the Godhead (Rom 2:16). Even the angels stand beneath His judiciary power in His sacred humanity, because His soul stands closest to the Word, His Passion exalted Him above every creature (Phil 2:10), and the angels minister to men whose Head He is.\endnote{Aquinas, \textit{ST} III, q.\ 59, a.\ 2, ad 3; a.\ 4, corpus and ad 1; a.\ 6, corpus.}

In what form will the Judge appear? St.\ Cyril of Jerusalem contrasts Christ's two advents:

\begin{quote}
We preach not one advent only of Christ, but a second also, far more glorious than the former. For the former gave a view of His patience; but the latter brings with it the crown of a divine kingdom. \dots\ In His former advent, He was wrapped in swaddling clothes in the manger; in His second, He covereth Himself with light as with a garment. In His first coming, He endured the Cross, despising shame; in His second, He comes attended by a host of Angels, receiving glory. \dots\ The Saviour comes, not to be judged again, but to judge them who judged Him.\endnote{St.\ Cyril of Jerusalem, \textit{Catechetical Lectures} XV.1, in \textit{NPNF} II, vol.\ 7.}
\end{quote}

Yet glory does not erase Calvary. The Apocalypse binds His majesty directly to His wounds: ``Behold, he cometh with the clouds, and every eye shall see him: and they also that pierced him. And all the tribes of the earth shall bewail themselves because of him'' (Apoc 1:7; cf.\ Zach 12:10; John 19:37). The \textit{Supplement} teaches that Christ will appear in the very flesh in which He was crucified, now radiant in glory and visible to all, bearing the five scars of His Passion:

\begin{quote}
The scars which will appear in His body will not be due to weakness, but will indicate the exceeding power whereby Christ overcame His enemies by His Passion and infirmity. He will also show forth His most shameful death, not by bringing it sensibly before the eye, as though He suffered it there; but by the things which will appear then, namely the signs of His past Passion, He will recall men to the thought of His past death.\endnote{\textit{Supplementum}, q.\ 90, a.\ 2, ad 2.}
\end{quote}

Those glorious wounds divide humanity before a word is spoken: ``as the sight of the glory of Christ's humanity will be a reward to the righteous, so will it be a torment to Christ's enemies.''\endnote{\textit{Supplementum}, q.\ 90, a.\ 2, ad 4.} To His friends, those scars are the living pledge of their ransom; to the impenitent, they are the unanswerable indictment of mercy spurned.
